\documentclass[../Note.tex]{subfiles}
\usepackage{bm}

\begin{document}

\chapter{Discrete Parameter Markov Chain} 

\begin{definition}[Stochastic Process]
    A stochastic process $\bm{X} = \{X(t), t \in \mathcal{T}\}$ is a collection of random variables. That is, for each $t$ in the index set $\mathcal{T}$, $X(t)$ is a random variable.
\end{definition}

The index is often interpreted as \textit{time} and, as a result, we refer to $X(t)$ as the \textit{state} of the process at time $t$. The \textit{state space} of a stochastic process is defined as the set of all possible values that the random variables $X(t)$ can assume. For instance,
\begin{itemize}
    \item $\{X(t), t = 0, 1, \dots\}$ is a discrete-time stochastic process indexed by the nonnegative integers;
    \item $\{X(t), t \geq 0\}$ is a continuous-time stochastic process indexed by the nonnegative real numbers. 
\end{itemize}
Thus, a stochastic process is a family of random variables that describes the system evolution through time of some (physical) process. This gives the conceptual fundation for \textit{Time-Dependent System Reliability}.

\textit{Markov Property} states that given present state, forecast of future is independent of past, i.e.,
\[
    \mathrm{P}\{\text{Futhure | \text{Present, Past}}\} = \mathrm{P}\{\text{Futhure | Present}\}.
\]
Then, a stochastic process that follows the Markov property is called a \textit{Markov process}.

\begin{definition}[Discrete-Parameter Markov Chain]
    Let $\{X_n, n = 0, 1, \dots\}$ be a stochastic process that takes on a finite or countable number of possible values. Suppose that 
    \[
        \mathrm{P}\{X_{n+1} = j | X_{n} = i, X_{n-1} = i_{n-1}, \dots, X_{1} = i_{1}, X_{0} = i_{0}\} = \mathrm{P}\{X_{n+1} = j | X_{n} = i\}
    \]
    for all states $i_0, i_i, \dots, i, j$ and all $n \geq 0$, then stochastic process $\{X_n, n = 0, 1, \dots\}$ is known as a discrete-parameter Markov chain (DPMC). In addition, if $\mathrm{P}\{X_{n+1} = j | X_{n} = i\} = p_{ij}$ for any $n \geq 0$, it is a homogeneous DPMC.
\end{definition}

In homogeneous DPMC, $p_{ij}$ represents the probability that the process will when in state $i$, next make a transition into state $j$, which is call \textit{one-step transition probability}. Let $\mathbf{P}$ denote the matrix of $p_{ij}$, so that the \textit{transition probability matrix} is 
\[
    P=
    \begin{bmatrix}
    p_{00} & p_{01} & p_{02} & \cdots \\
    p_{10} & p_{11} & p_{12} & \cdots \\
    \vdots & \vdots & \vdots & \vdots \\
    p_{i0} & p_{i1} & p_{i2} & \cdots \\
    \vdots & \vdots & \ddots & \vdots
    \end{bmatrix},
\]
where $0 \leq p_{ij} \leq 1$ and $\sum_{j} p_{ij} = 1$ for all $i$.

Define the $n$-step transition probability $p_{ij}(n)$ as 
\[
    p_{ij}(n) = \mathrm{P}\{X_{n+k} = j | X_k = i\}, \quad n \geq 0, i,j \geq 0.
\]

\begin{theorem}[Chapman-Kolmogorov Equations]
    The \textit{Chapman-Kolmogorov Equations} are 
    \begin{equation}
        p_{ij}(n+m) = \sum\limits_{k=0}^{\infty} p_{ik}{(n)}p_{kj}{(m)} \quad \forall n,m \ge 0, \forall i,j.
    \end{equation} 
\end{theorem}

\begin{proof}
    \[
        \begin{aligned}
            p_{ij}(n+m) = &  \mathrm{P}\{X_{n+m} = j | X_0 = i\} \qquad && \leftarrow \text{Definition} \\
            & = \sum_{k = 0}^{\infty} \mathrm{P}\{X_{n+m} = j, X_n = k | X_0 = i\} \qquad && \leftarrow \text{Law of total probability} \\
            & = \sum_{k = 0}^{\infty} \mathrm{P}\{X_{n+m} = j | X_n = k, X_0 = i\} \mathrm{P}\{X_n = k | X_0 = i\} \qquad && \leftarrow \text{Conditional probability}  \\
            & = \sum_{k = 0}^{\infty} \mathrm{P}\{X_{n+m} = j | X_n = k\} \mathrm{P}\{X_n = k | X_0 = i\} \qquad && \leftarrow \text{Markov property} \\
            & = \sum\limits_{k=0}^{\infty} p_{ik}{(n)}p_{kj}{(m)}.
        \end{aligned}    
    \] 
\end{proof}

If we let $\mathbf{P}(n)$ denote the matrix of $n$-step transition probability $p_{ij}(n+m)$, then 
\begin{equation}
    \mathbf{P}(n+m) = \mathbf{P}(n) \cdot \mathbf{P}(m),
\end{equation} 
in particular,
\[
    \mathbf{P}(2) = \mathbf{P}^{2} \qquad \mathbf{P}(n) = \mathbf{P}^{n}.
\]

\begin{remark}
    The matrix of $n$-step transition probabilities is obtained by multiplying the matrix of one-step transition probabilities by itself $n-1$ times.
\end{remark}

Let $p_i(n)$ denote the probability that a Markov chain is in state $i$ at time $m$, i.e.,
\[
    p_i(n) = \mathrm{P} \left\{X_n = i\right\}.
\]
In vector notation, let $\bm{p}(n) = (p_1(n), p_2(n), \dots, p_i(n), \dots)$. 

The state probabilities at any time step $n$ can be obtained from a knowledge of the initial probability distribution and the matrix of transition probabilities. From the theorem of total probability,
\[
    p_i(n) = \sum_{\text{all } k} p_k(0) \mathrm{P}\{X_n = i | X_0 = k\} = \sum_{\text{all } k} p_k(0) p_{ki}(n),
\]
writing in matrix form becomes
\[
    \bm{p}(n) = \bm{p}(0) \mathbf{P}(n) = \bm{p}(0) \mathbf{P}^{n},
\]
where $\bm{p}(0)$ denotes the initial state distribution.

\begin{definition}[Steady-State Distribution]
    Let $\mathbf{P}$ be the transition probability matrix of a DPMC, and let the vector $\bm{\pi}$ be a probability distribution, whose element $\pi_j$ denote the probability of being in state $j$, i.e., 
    \[
        0 \leq \pi_j \leq 1, \qquad \sum_{\text{all } j } \pi_j = 1.
    \]
    Then $\bm{\pi}$ is said to be steady-state distribution if and only if 
    \[
        \bm{\pi} \mathbf{P} = \bm{\pi}.
    \]
\end{definition} 

\begin{remark}[Steady-State Distribution \& Eigenproblem]
    Recall the eigenvalue problem
    \[
        \mathbf{A}\bm{x}=\lambda\bm{x},
    \] 
    where $\lambda$ is an eigenvalue of $\mathbf{A}$ and $\bm{x}$ is the corresponding right eigenvector.

    Since a steady-state distribution satisfies
    \[
        \bm{\pi}\mathbf{P}=\bm{\pi},
    \] 
    $\bm{\pi}$ is a left eigenvector of $\mathbf{P}$ associated with the eigenvalue $\lambda=1$.

    Equivalently,
    \[
        \mathbf{P}^{\top}\bm{\pi}^{\top} = \bm{\pi}^{\top}.
    \]
    Thus, the stationary distribution can be computed by solving
    \begin{equation}
    \left\{
    \begin{aligned}
        \bm{\pi}\mathbf{P} &= \bm{\pi}, \\
        \bm{\pi}\bm{e} &= 1.
    \end{aligned}
    \right.        
    \end{equation} 
    where $\bm{e}$ is the vector of ones.
\end{remark}

Define $f_{ij}(n)$ for $i \ne j$ as the probability that, starting from state $i$, the first passage to state $j$ occurs in exactly $n$ steps 
\[ 
    f_{ij}(n) = P\{X_n = j , X_k \ne j , k = 1, \dots, n -1 | X_0 = i\}, \qquad n = 1, 2, \dots.
\]

\begin{proposition}[Relationship between $f_{ij}(n)$ and $p_{ij}(n)$]
    \label{proposition-1-1}
    For $i \ne j$,
    \begin{equation}
        \label{eq-01-04}
        f_{ij}(n) = p_{ij}(n) - \sum_{l = 1}^{n-1} f_{ij}(l) p_{ij}(n-l).
    \end{equation}
    When $i=j$,
    \begin{equation}
        f_{ii}(n) = p_{ii}(n) - \sum_{l = 1}^{n-1} f_{ii}(l) p_{ii}(n-l).
    \end{equation}
\end{proposition} 
\begin{proof}
    If the chain is in state $j$ at step $n$, then the first visit to state $j$ must occur at some step $\ell \le n$. If the first visit occurs at step $\ell$, then the probability of this event is $f_{ij}(\ell)$. After reaching state $j$, the probability of being in state $j$ again after the remaining $n-\ell$ steps is
    $p_{jj}(n-\ell)$. Therefore,
    \[
        \begin{aligned}
            p_{ij}(1)
            &= f_{ij}(1), \\
            p_{ij}(2)
            &= f_{ij}(1)p_{jj}(1) + f_{ij}(2), \\
            p_{ij}(3)
            &= f_{ij}(1)p_{jj}(2)
            + f_{ij}(2)p_{jj}(1)
            + f_{ij}(3), \\
            &\vdots \\
            p_{ij}(n)
            &= f_{ij}(n)
            + \sum_{\ell=1}^{n-1}
            f_{ij}(\ell)p_{jj}(n-\ell).
        \end{aligned}
    \]
    Rearranging gives
    \[
        f_{ij}(n)
        =
        p_{ij}(n)
        -
        \sum_{\ell=1}^{n-1}
        f_{ij}(\ell)p_{jj}(n-\ell).
    \]

    The case $i=j$ follows analogously.
\end{proof}

Proposition \ref{proposition-1-1} essentially gives the following relationship:
\[
\begin{aligned}
    \text{Being in state } j \text{ at step } n = \sum_{\ell=1}^{n} \left( \text{first reaching } j \text{ at step } \ell \right) \times
    \left( \text{being in } j \text{ again after the remaining } n-\ell \text{ steps} \right),
\end{aligned}
\]
and it also gives a recursive way to compute $f_{ij}(n)$ based on $p_{ij}(n)$.

\begin{remark}[How to Calculate $p_{ij}(n)$ and $f_{ij}(n)$]
    Use CK equation in matrix form for $p_{ij}(n)$, and then use the recursive equation \eqref{eq-01-04} for $f_{ij}(n)$.
\end{remark}

Let 
\[
    T_j = \min\{n \geq 1: X_n = j\}
\]
be the first return time to $j$. Then, we can define the mean recurrence time $m_{jj}$ of state $j$ as 
\[
    m_{jj} = \mathbb{E}[\text{Recurrence time}] = \mathbb{E}[T_j] = \sum_{n = 1}^{\infty} n f_{jj}(T_j = n) = \sum_{n = 1}^{\infty} n f_{jj}(n),
\]
which is the average number of steps taken to return to state $j$ for the first time after leaving it. Similarly, for $i \ne j$, we can define the mean first passage time
\[
    m_{ij} = \sum_{n = 1}^{\infty} n f_{ij}(n).
\]
The $m_{ij}$ uniquely satisfy
\[
    m_{ij} = f_{ij}(1) + \sum_{k \ne j} p_{ik} (1 + m_{kj}) = 1 + \sum_{k \ne j}p_{ik}m_{kj},
\]
since the process in state $i$ either goes to state $j$ in one step with probability $p_{ij}$, or goes first to some intermediate state $k$ in one step with probability $p_{ik}$ and then eventually on to $j$ in an additional $m_{kj}$ steps.

Define $q_{ij}(m)$ as the probability that, starting from state $i$, state $j$ is reached for the first time within $m$ steps
\[
    q_{ij}(m) = \sum_{n=1}^{m} f_{ij}(n).
\]
Equivalently,
\[
    q_{ij}(m) = P\{T_j \leq m \mid X_0=i\},
\]
where $T_j$ denotes the first passage time to state $j$. Since $q_{ij}(m)$ is the cumulative probability associated with the first passage probabilities $f_{ij}(n)$,
\[
    f_{ij}(m) = q_{ij}(m)-q_{ij}(m-1).
\]

Letting $m\to\infty$, we define
\[
    f_{ij}=\sum_{n=1}^{\infty} f_{ij}(n)=\lim_{m\to\infty} q_{ij}(m),
\]
which is the probability that state $j$ is ever reached from state $i$. When $i=j$,
\[
    f_{jj}=\sum_{n=1}^{\infty} f_{jj}(n)=\lim_{m\to\infty}q_{jj}(m),
\]
which is the probability that, starting from state $j$, the process eventually returns to state $j$.

\begin{proposition}[Relationship among $p_{ij}(n)$, $f_{ij}(n)$, and $q_{ij}(n)$]
    For $i \ne j$ and $n \geq 1$
    \[
        q_{ij}(n) \geq p_{ij}(n) \geq f_{ij}(n).
    \]
\end{proposition}
\begin{proof}
    We prove this directly from the definitions. Recall that
    \[
        \begin{aligned}
            p_{ij}(n)
            &= P(X_n = j \mid X_0 = i), \\
            f_{ij}(n)
            &= P(X_n = j,\ X_k \neq j,\ k=1,\dots,n-1 \mid X_0 = i), \\
            q_{ij}(n)
            &= P\left(\min\{m : X_m = j\} \leq n \mid X_0 = i\right).
        \end{aligned}
    \]

    First, split $q_{ij}(n)$ according to whether the first visit to $j$
    occurs before time $n$ or exactly at time $n$:
    \[ 
        q_{ij}(n) = P\left(\min\{m : X_m=j\}\leq n-1 \mid X_0=i\right) + \underbrace{P\left(\min\{m : X_m=j\}=n \mid X_0=i\right)}_{f_{ij}(n)}. 
    \]

    Next, split $p_{ij}(n)$ according to whether the chain visits $j$
    for the first time at time $n$, or has already visited $j$ by time
    $n-1$:
    \[ 
        p_{ij}(n) = P(X_n=j \mid X_0=i) = \underbrace{P\left(\min\{m : X_m=j\}=n \mid X_0=i\right)}_{f_{ij}(n)} + P\left(\min\{m : X_m=j\}\leq n-1,\ X_n=j \mid X_0=i\right). 
    \]

    Since
    \[
        \left\{
            \min\{m : X_m=j\}\leq n-1,\ X_n=j
        \right\}
        \subseteq
        \left\{
            \min\{m : X_m=j\}\leq n-1
        \right\},
    \]
    we have
    \[
        P\left(
            \min\{m : X_m=j\}\leq n-1,\ X_n=j
            \mid X_0=i
        \right)
        \leq
        P\left(
            \min\{m : X_m=j\}\leq n-1
            \mid X_0=i
        \right).
    \]

    Therefore, comparing the two decompositions above gives
    \[
        p_{ij}(n) \leq q_{ij}(n).
    \]

    Also, from
    \[
        p_{ij}(n)
        =
        f_{ij}(n)
        +
        P\left(
            \min\{m : X_m=j\}\leq n-1,\ X_n=j
            \mid X_0=i
        \right),
    \]
    and the nonnegativity of probability, we obtain
    \[
        p_{ij}(n) \geq f_{ij}(n).
    \]

    Hence,
    \[
        q_{ij}(n) \geq p_{ij}(n) \geq f_{ij}(n).
    \]
\end{proof}

Let 
\[
    L_j = \min \{n \geq 1 : X_n \ne j\}
\]
be starting from state $j$, the number of steps until the process first leaves state $j$. Then, we can define the mean occupation time $l_{jj}$ of state $j$ as
\[ 
    l_{jj} = \mathbb{E}[\text{Occupation time}] = \mathbb{E}[L_j] = \sum_{n = 1}^{\infty} n p_{jj}^{n-1}(1-p_{jj}) = \frac{1}{1-p_{ii}}. 
\]

\begin{proposition}[Closed-Form Expression of Mean Occupation Time]
    The closed-form of mean occupation time of state $j$ is
    \[
        l_{jj}=\mathbb{E}[L_j]=\frac{1}{1-p_{jj}}.
    \]
\end{proposition}

\begin{proof}
    Note that $L_{j} \sim \mathrm{Geom}(1 - p_{jj})$ and 
    \[
        \sum_{n=0}^{\infty}x^n = \frac{1}{1-x}, \quad |x|<1.
    \]
    Differentiate both sides with respect to $x$ to obtain
    \[
        \sum_{n=1}^{\infty} n x^{n-1}=\frac{1}{(1-x)^2}.
    \]
    Letting $x=p_{jj}$,
    \[
        \sum_{n=1}^{\infty}n p_{jj}^{\,n-1}=\frac{1}{(1-p_{jj})^2}.
    \]
    Therefore,
    \[ 
        l_{jj}=\sum_{n=1}^{\infty}n p_{jj}^{\,n-1}(1-p_{jj}) =(1-p_{jj})\sum_{n=1}^{\infty}n p_{jj}^{n-1} =(1-p_{jj})\frac{1}{(1-p_{jj})^2} =\frac{1}{1-p_{jj}}. 
    \]
\end{proof}




\end{document}